% Table 3: H0 Measurement Compilation and Multi-Method Convergence
% Generated from h0_measurements_compilation.csv

\begin{deluxetable*}{lccc}
\tablecaption{H$_0$ Measurement Compilation and Multi-Method Convergence\label{tab:h0_compilation}}
\tablewidth{0pt}
\tablehead{
\colhead{Method} &
\colhead{H$_0$} &
\colhead{$\sigma$} &
\colhead{Reference} \\
\colhead{} &
\colhead{(km~s$^{-1}$~Mpc$^{-1}$)} &
\colhead{(km~s$^{-1}$~Mpc$^{-1}$)} &
\colhead{}
}
\startdata
SH0ES Cepheid & 73.04 & 1.04 & \citet{Riess2022} \\
TRGB & 69.85 & 2.33 & \citet{Freedman2025a} \\
JAGB & 67.96 & 2.65 & \citet{Freedman2025a} \\
Cosmic Chronometers (H(z)) & 68.33 & 1.57 & This work \\
Planck CMB & 67.36 & 0.54 & \citet{Planck2018} \\
\hline
\multicolumn{4}{c}{\textit{Late-Universe (Planck-Independent)}} \\
\hline
JAGB + Cosmic Chron. & 68.22 & 1.36 & This work \\
\hline
\multicolumn{4}{c}{\textit{Three-Method Convergence (incl. Planck)}} \\
\hline
Weighted Mean & 67.48 & 0.50 & This work \\
\enddata
\tablecomments{Compilation of H\ensuremath{_0} measurements from different methods revealing systematic gradient and convergence. SH0ES Cepheid is shown with the published Riess et al.\ (2022) uncertainty (1.04\,km\,s\ensuremath{^{-1}}\,Mpc\ensuremath{^{-1}}). Corrected Cepheid (baseline Scenario A + Prior 1, not shown in the main rows) is \ensuremath{H_0 = 69.54 \pm 1.89}\,km\,s\ensuremath{^{-1}}\,Mpc\ensuremath{^{-1}} after applying realistic correlated systematics and the period and metallicity bias corrections. JAGB + cosmic chronometers give a Planck-independent mean of \ensuremath{68.22 \pm 1.36}\,km\,s\ensuremath{^{-1}}\,Mpc\ensuremath{^{-1}} with \ensuremath{\chi^2_{\rm red} \approx 0.01}; the corrected Cepheid value differs from this Planck-free mean by \ensuremath{0.6\sigma}. The three-method weighted mean of JAGB, cosmic chronometers, and Planck is \ensuremath{67.48 \pm 0.50}\,km\,s\ensuremath{^{-1}}\,Mpc\ensuremath{^{-1}} with \ensuremath{\chi^2_{\rm red} = 0.19}, showing strong cross-method convergence.}
\end{deluxetable*}
